%% ============================================================================
%%  Cost Function Design for Bistatic Receiver Placement in Helical SAR
%%  Subsurface Imaging Systems
%%
%%  IEEE Transactions on Geoscience and Remote Sensing — manuscript format
%%  Author: David Atencia
%% ============================================================================
\documentclass[journal]{IEEEtran}

%% ─── Encoding ───────────────────────────────────────────────────────────────
\usepackage[utf8]{inputenc}
\usepackage[T1]{fontenc}

%% ─── Mathematics ────────────────────────────────────────────────────────────
\usepackage{amsmath}
\usepackage{amssymb}
\usepackage{mathtools}
\usepackage{bm}

%% ─── Figures and tables ─────────────────────────────────────────────────────
\usepackage{graphicx}
\usepackage{booktabs}
\usepackage{array}
\usepackage{caption}
\usepackage{subcaption}
\usepackage{xcolor}
\usepackage{float}

%% ─── Hyperlinks ─────────────────────────────────────────────────────────────
\usepackage[hidelinks, colorlinks=true, linkcolor=blue!60!black,
            citecolor=blue!60!black, urlcolor=blue!60!black]{hyperref}

%% ─── Macros ─────────────────────────────────────────────────────────────────
\newcommand{\Rx}{\mathrm{Rx}}
\newcommand{\Tx}{\mathrm{Tx}}
\newcommand{\rRx}{\mathbf{r}_{\Rx}}
\newcommand{\rTx}{\mathbf{r}_{\Tx}}
\newcommand{\dxy}{\delta_{xy}}
\newcommand{\dz}{\delta_z}
\newcommand{\gdeg}{\ensuremath{^{\circ}}}
\newcommand{\thetaB}{\theta_{\mathrm{B}}}
\newcommand{\norm}[1]{\left\|#1\right\|}
\DeclareMathOperator*{\argmin}{arg\,min}

%% ─── Metadata ───────────────────────────────────────────────────────────────
\title{Cost Function Design for Bistatic Receiver Placement\\
       in Helical SAR Subsurface Imaging Systems}

\author{David~Atencia%
  \thanks{Manuscript received \today.}%
  \thanks{The author is with the Department of Electrical Engineering,
          [Institution], [City], [Country]
          (e-mail: david.atencia.m21@gmail.com).}}

\markboth{IEEE Transactions on Geoscience and Remote Sensing}%
         {Atencia: Cost Function Design for Bistatic SAR Receiver Placement}

% ============================================================================
\begin{document}
% ============================================================================

\maketitle

% ─── Abstract ────────────────────────────────────────────────────────────────
\begin{abstract}
This paper investigates the design of scalarization cost functions for the
problem of optimal bistatic receiver placement in helical synthetic aperture
radar (SAR) systems intended for subsurface imaging. The system geometry
consists of a fixed transmitter following a circular orbit and a configurable
airborne receiver whose horizontal position simultaneously determines two
conflicting imaging quality objectives: the received signal power, governed by
the TM-polarization Brewster transmittance at the air--soil interface, and the
spatial resolution, controlled by the bistatic angle and the elevation geometry.
Five candidate cost functions are proposed and analyzed, spanning weighted-sum
scalarization in logarithmic and linear-normalized scales, Chebyshev minimax
scalarization, ideal-point Euclidean distance, and a Brewster-penalized
augmented form. For each function, the optimization landscape, Pareto-front
coverage, sensitivity to the trade-off parameter~$\alpha$, and gradient
conditioning are evaluated through analytical simulation over a representative
airborne bistatic SAR geometry at X-band (10~GHz) with a dry-soil half-space
model ($\varepsilon_r = 4$). Results show that the logarithmic weighted-sum
cost function ($J_1$) provides the best balance of differentiability, intrinsic
commensurability across orders-of-magnitude variations in received power, and
unimodal landscape topology, making it the recommended choice for
gradient-based optimization under the two-layer propagation model. The
Chebyshev scalarization is identified as the preferred alternative for
multilayer geometries exhibiting non-convex Pareto fronts.
\end{abstract}

% ─── Index Terms ─────────────────────────────────────────────────────────────
\begin{IEEEkeywords}
Bistatic SAR, synthetic aperture radar, ground-penetrating radar, subsurface
imaging, multiobjective optimization, Pareto front, Brewster angle, TM
polarization, receiver placement, cost function, Chebyshev scalarization.
\end{IEEEkeywords}

% ============================================================================
\section{Introduction}
\label{sec:intro}
% ============================================================================

\IEEEPARstart{S}{ynthetic} aperture radar (SAR) has established itself as the
principal remote-sensing modality for large-area subsurface imaging, combining
the all-weather, day-and-night operational capability of microwave systems with
the spatial resolution achievable through synthetic aperture
processing~\cite{moreira2013, cumming2005}. Applications range from
archaeological prospection and planetary exploration to security screening and
unexploded-ordnance detection, each imposing demanding requirements on the
signal-to-noise ratio and the volumetric spatial resolution of the
imaged scene~\cite{daniels2004, yarovoy2007}.

Conventional monostatic ground-penetrating radar (GPR) systems transmit and
receive from the same platform, limiting the angular diversity of the
measurement. The bistatic configuration—where the transmitter (Tx) and
receiver (Rx) are physically separated—offers substantially higher degrees of
freedom: the bistatic angle introduces additional wavevector diversity that
directly improves cross-range resolution~\cite{willis1991, krieger2006}.
Furthermore, when the medium is lossy or the ground surface is rough, separating
Tx and Rx allows the system to exploit the angular dependence of the Fresnel
transmission coefficient, and in particular the TM-polarization Brewster
condition, at which the air--soil interface becomes perfectly transparent and
the received signal power is maximized~\cite{born1999, balanis2012}.

The helical bistatic SAR architecture studied here places the Tx on a fixed
circular orbit while the Rx occupies a configurable position at the same
altitude. The orbital aperture, combined with a controlled receiver offset,
generates the synthetic aperture required for three-dimensional reconstruction.
However, the position of the Rx is the key free variable of the design: it
simultaneously controls the received signal power (through the Brewster
transmittance) and the spatial resolution (through the bistatic angle and
elevation geometry). These two objectives are fundamentally conflicting, because
the Brewster condition occurs at a specific elevation angle that does not
generally coincide with the geometry maximizing bistatic resolution.

Optimization of bistatic radar geometries has been studied in the context of
target detection~\cite{willis1991}, interferometric baseline
design~\cite{krieger2006}, and tomographic aperture synthesis~\cite{reigber2000,
tebaldini2012}. However, the \emph{joint} optimization of power and resolution
under realistic refraction constraints—solved as a multi-objective problem with
an explicit cost function landscape analysis—has not been comprehensively
addressed in the literature.

Multi-objective optimization problems can be approached through scalarization,
converting the vector objective into a single scalar that is then minimized by a
standard optimizer. The choice of scalarization function is non-neutral: it
determines the subset of the Pareto front that is reachable, the topology of
the optimization landscape, the gradient conditioning for local-search methods,
and the physical interpretability of the trade-off parameter. A systematic
comparison of scalarization strategies tailored to a physically motivated
bistatic SAR geometry represents the principal contribution of this work.

The contributions of this paper are as follows:
\begin{enumerate}
  \item Formal definition of the bistatic receiver placement problem as a
    two-objective optimization in which signal power and volumetric resolution
    are expressed as commensurable logarithmic scalars.
  \item Proposal and theoretical analysis of five candidate cost functions
    spanning the main scalarization families used in the multiobjective
    optimization literature.
  \item Comparative simulation study of the optimization landscapes, Pareto
    fronts, $\alpha$-sensitivity curves, and gradient norms for each function
    under a representative X-band helical bistatic SAR geometry.
  \item Practical recommendations for cost function selection, including a
    physically motivated alternative ($\varepsilon$-constrained formulation)
    and extension guidelines for multilayer subsurface models.
\end{enumerate}

The remainder of the paper is organized as follows. Section~\ref{sec:model}
describes the system geometry and physical models. Section~\ref{sec:problem}
formulates the multi-objective optimization problem. Section~\ref{sec:cost}
defines and analyzes the five candidate cost functions.
Section~\ref{sec:results} presents the simulation results. Section~\ref{sec:discussion}
discusses the findings and provides design recommendations. Section~\ref{sec:conclusion}
concludes the paper.

% ============================================================================
\section{System Model}
\label{sec:model}
% ============================================================================

\subsection{Helical Bistatic SAR Geometry}

The system consists of two airborne platforms, both flying at altitude
$z_{\Tx} = z_{\Rx} = H = 100\,\mathrm{m}$ above a flat air--soil interface
located at $z = 0$. The transmitter follows a horizontal circular orbit of
radius $\rho_{\Tx} = 160\,\mathrm{m}$ centered over a subsurface target of
interest located at $\mathbf{r}_P = (0, 0, -z_P)$ with $z_P = 5\,\mathrm{m}$.
The orbital motion provides the coherent aperture necessary for SAR processing.
The receiver position $\rRx = (x_{\Rx}, y_{\Rx}, H)$ is a two-dimensional
design parameter with $(x_{\Rx}, y_{\Rx}) \in [-270, 270]^2\,\mathrm{m}$.

The combined synthetic aperture formed by the Tx orbit and the fixed Rx
position generates a \emph{helical} aperture in the bistatic wavenumber space,
whose shape determines the three-dimensional point spread function (PSF) of the
reconstructed image~\cite{soumekh1999}. The bistatic azimuth angle
$\Delta\phi = \phi_{\Rx} - \phi_{\Tx}$, measured from the target, is the
principal geometric variable coupling power and resolution.

\subsection{TM Transmittance and Brewster Condition}

The soil is modeled as a homogeneous half-space with complex permittivity
$\varepsilon_r = n_2^2 = 4$ (dry mineral soil at X-band~\cite{ulaby2014}),
giving a refractive index $n_2 = 2$ and a Brewster angle

\begin{equation}
  \thetaB = \arctan\!\left(\frac{n_2}{n_1}\right) = 63.43\gdeg
  \label{eq:thetaB}
\end{equation}

where $n_1 = 1$ for air. At TM polarization, the Fresnel transmission
coefficient is

\begin{equation}
  t_{\mathrm{TM}}(\theta_i) =
    \frac{2\,n_1\cos\theta_i}
         {n_2\cos\theta_i + n_1\cos\theta_t}
  \label{eq:tTM}
\end{equation}

with the refracted angle $\theta_t = \arcsin(n_1\sin\theta_i/n_2)$ given by
Snell's law. The power transmittance $T = |t_{\mathrm{TM}}|^2 \cos\theta_t /
(n_1 \cos\theta_i / n_2)$ reaches unity at $\theta_i = \thetaB$~\cite{born1999},
making this angle critical for maximizing penetration depth and received power.

The ray path from the Rx to the subsurface target is refracted at the
air--soil interface. The refracted ray parameter $p = n_1\sin\theta_i$ is
found by solving Fermat's principle

\begin{equation}
  \frac{\partial}{\partial p}\!\left[
    \frac{\sqrt{\rho_{\mathrm{hs}}^2 + H^2}}{\,n_1/p\,}
    + \frac{\sqrt{(n_2/p)^2 - 1}\cdot z_P}{n_2}
  \right] = \rho_{\Rx P}
  \label{eq:fermat}
\end{equation}

numerically via bisection, where $\rho_{\mathrm{hs}}$ is the horizontal
distance from Rx to the surface reflection point and $\rho_{\Rx P}$ is the
horizontal distance from Rx to the target projected on the surface.

\subsection{Bistatic Radar Equation}

The received power at $\rRx$ for a point target of radar cross-section
$\sigma$ follows the bistatic radar equation modified for two-way refraction:

\begin{equation}
  P_r(\rRx) = \frac{P_t\,G_t\,G_r\,\sigma\,T_1\,T_2(\rRx)\,\lambda_0^2}
                   {(4\pi)^3\,R_T^2\,R_R^2(\rRx)}
  \label{eq:Pr}
\end{equation}

where $P_t$ is the transmitted power, $G_t$ and $G_r$ are antenna gains,
$\lambda_0 = c/f_0$ is the free-space wavelength, $T_1$ is the TM
transmittance on the Tx$\to$target path (constant for fixed Tx), $T_2(\rRx)$
is the transmittance on the Rx$\to$target path, and $R_T$, $R_R(\rRx)$ are
the respective slant-range path lengths computed along the refracted ray.

\subsection{Spatial Resolution}

The volumetric spatial resolution is derived from the $k$-space support of
the bistatic measurement operator~\cite{soumekh1999, devaney1984}. The axial
(depth) and lateral resolutions are, respectively:

\begin{align}
  \dz &= \frac{c}{2\,W_z}, \quad
  W_z = n_2\,B\cos\theta_{t,0}
        + \frac{c\,B_{\mathrm{hel}}\sin\psi_0\cos\psi_0}
               {\lambda_0\,R_0\,n_2\cos\theta_{t,0}}
  \label{eq:dz} \\[4pt]
  \dxy &= \frac{0.60\,\lambda_0}
               {\pi\sin\psi_0 \cdot |\cos(\Delta\phi/2)|}
  \label{eq:dxy}
\end{align}

where $\psi_0$ is the elevation angle from the Rx to the surface reflection
point, $\theta_{t,0} = \arcsin(\sin\psi_0/n_2)$ is the refracted angle at the
interface, $R_0$ is the oblique range from the Rx to the surface, and
$B_{\mathrm{hel}} = 47.17\,\mathrm{m}$ is the effective helical aperture.

The lateral resolution $\dxy$ improves as $|\cos(\Delta\phi/2)|$ increases,
i.e., as the bistatic angle $\Delta\phi$ approaches $180\gdeg$. The axial
resolution $\dz$ is primarily controlled by the bandwidth $B$ and the
refraction geometry at the Rx elevation.

A summary of the system reference parameters is provided in
Table~\ref{tab:params}.

\begin{table}[!t]
  \renewcommand{\arraystretch}{1.2}
  \centering
  \caption{Reference System Parameters}
  \label{tab:params}
  \begin{tabular}{@{}lll@{}}
    \toprule
    \textbf{Parameter} & \textbf{Value} & \textbf{Description} \\
    \midrule
    $f_0$                  & 10\,GHz       & Carrier frequency \\
    $B$                    & 50\,MHz       & Chirp bandwidth \\
    $n_1$                  & 1.0           & Refractive index, air \\
    $n_2$                  & 2.0           & Refractive index, soil \\
    $\varepsilon_r$        & 4.0           & Relative permittivity \\
    $\thetaB$              & 63.43\gdeg    & TM Brewster angle \\
    $z_{\Tx} = z_{\Rx}$   & 100\,m        & Platform altitude \\
    $\rho_{\Tx}$           & 160\,m        & Tx orbit radius \\
    $B_{\mathrm{hel}}$     & 47.17\,m      & Helical aperture \\
    $z_P$                  & 5\,m          & Target depth \\
    \bottomrule
  \end{tabular}
\end{table}

% ============================================================================
\section{Problem Formulation}
\label{sec:problem}
% ============================================================================

\subsection{Elementary Objective Functions}

The receiver placement problem is formulated as a biobjective optimization in
which each objective is expressed as a logarithmic scalar to achieve
dimensional commensurability across the large dynamic range of $P_r$.

\subsubsection{Power-loss objective}

\begin{equation}
  f_1(\rRx) = -\log_{10} P_r(\rRx) \geq 0
  \label{eq:f1}
\end{equation}

Minimizing $f_1$ is equivalent to maximizing the received power. The
logarithmic mapping compresses the several orders of magnitude spanned by
$P_r$ over the search domain (typically $10^{-17}$ to $10^{-14}$~W for the
reference parameters) into the bounded interval $f_1 \in [14, 18]$.

\subsubsection{Resolution-volume objective}

\begin{equation}
  f_2(\rRx) = \log_{10}\!\bigl[\dxy(\rRx) \cdot \dz(\rRx)\bigr]
  \label{eq:f2}
\end{equation}

Minimizing $f_2$ is equivalent to minimizing the resolution cell volume.
For the reference system, $f_2 \in [-3, -1.5]$, providing magnitudes
comparable to $f_1$ without explicit normalization.

\subsection{Multi-Objective Formulation and Pareto Optimality}

The vector optimization problem is

\begin{equation}
  \min_{\rRx \in \mathcal{D}} \; \mathbf{f}(\rRx)
  = \bigl[f_1(\rRx),\; f_2(\rRx)\bigr]^{\top}
  \label{eq:mop}
\end{equation}

where $\mathcal{D} = [-270,270]^2\,\mathrm{m}$ is the search domain at fixed
altitude $H$. A position $\rRx^*$ is \emph{Pareto optimal} (non-dominated) if
no other feasible point achieves a strictly better value in both objectives
simultaneously~\cite{miettinen1998}. The set of all Pareto-optimal points forms
the \emph{Pareto front} $\mathcal{P}^*$ in the objective space.

The conflict between $f_1$ and $f_2$ arises because:
\begin{itemize}
  \item Maximum power (minimum $f_1$) is achieved when $\theta_R = \thetaB$,
    constraining the Rx to an annular region of radius
    $\rho_B \approx H\tan\thetaB \approx 200\,\mathrm{m}$ around the target.
  \item Minimum resolution volume (minimum $f_2$) requires large $\Delta\phi$
    (bistatic diversity), which constrains the Rx to be near-opposite to the
    Tx in azimuth, regardless of radial distance.
\end{itemize}

Solving~\eqref{eq:mop} directly requires mapping the full Pareto front, which
is computationally feasible via grid evaluation but unsuitable for embedded
real-time optimization. Scalarization converts~\eqref{eq:mop} into a
parametric single-objective problem that can be solved by any standard
optimizer.

% ============================================================================
\section{Candidate Cost Functions}
\label{sec:cost}
% ============================================================================

Five scalarization strategies are proposed and analyzed. Each represents a
well-studied family in the multiobjective optimization
literature~\cite{miettinen1998, deb2001, das1997}.

\subsection{\texorpdfstring{$J_1$}{J1} — Logarithmic Weighted Sum (Reference)}

\begin{equation}
  \boxed{J_1(\rRx;\,\alpha) = (1-\alpha)\,f_1 + \alpha\,f_2}
  \label{eq:J1}
\end{equation}

$J_1$ implements the classical weighted-sum scalarization~\cite{zadeh1963}
directly on the logarithmic objectives. The weight $\alpha \in [0,1]$ encodes
the designer's preference for resolution over power.

\subsubsection*{Properties}
\begin{itemize}
  \item \emph{Differentiability}: $J_1$ is differentiable wherever $P_r > 0$
    and $\dxy, \dz > 0$, i.e., over the entire valid search domain.
  \item \emph{Intrinsic commensurability}: The logarithmic scale ensures
    $f_1 \approx [14,18]$ and $f_2 \approx [-3,-1.5]$ are of comparable
    magnitude, eliminating the need for explicit normalization.
  \item \emph{Limiting cases}: $\alpha = 0$ yields the power-only solution
    (Brewster ring); $\alpha = 1$ yields the resolution-only solution
    ($\Delta\phi \to 180\gdeg$, degenerate case where the Rx moves to the
    domain boundary).
  \item \emph{Recommended value}: $\alpha^* \approx 0.30$ provides a balanced
    trade-off in the region of highest Pareto-front curvature.
  \item \emph{Limitation}: The weighted-sum method can only reach Pareto
    optimal points on convex portions of the Pareto front; non-convex segments
    are inaccessible for any $\alpha$~\cite{das1997}.
\end{itemize}

\subsection{\texorpdfstring{$J_2$}{J2} — Linear Normalized Weighted Sum}

\begin{equation}
  J_2(\rRx;\,\alpha) = (1-\alpha)\,\tilde{f}_1 + \alpha\,\tilde{f}_2
  \label{eq:J2}
\end{equation}

where $\tilde{f}_k = (f_k - f_k^{\min}) / (f_k^{\max} - f_k^{\min})$ is the
min-max normalization over~$\mathcal{D}$.

\subsubsection*{Properties}
\begin{itemize}
  \item \emph{Guaranteed commensurability}: both objectives mapped to $[0,1]$.
  \item \emph{Differentiability}: the normalization is an affine transformation,
    preserving differentiability.
  \item \emph{Critical limitation}: The min-max normalization depends on the
    search domain~$\mathcal{D}$. Changing the grid resolution or boundaries
    changes the cost function even though the physics does not, violating the
    optimizer invariance principle~\cite{nocedal2006}.
  \item \emph{Plateau regions}: In linear scale, $P_r$ varies over three orders
    of magnitude. After normalization, the vast majority of the domain maps to
    $\tilde{f}_1 \approx 0$, with only the Brewster neighborhood exhibiting
    significant gradient. This creates large flat regions (\emph{plateaus}) that
    degrade convergence of first-order methods~\cite{nocedal2006}.
\end{itemize}

\subsection{\texorpdfstring{$J_3$}{J3} — Weighted Chebyshev Scalarization (Minimax)}

\begin{equation}
  J_3(\rRx;\,\alpha) =
    \max\!\bigl[(1-\alpha)\,\tilde{f}_1,\;\alpha\,\tilde{f}_2\bigr]
  \label{eq:J3}
\end{equation}

with the same min-max normalization as $J_2$. This is the Chebyshev
($\ell_\infty$) scalarization~\cite{messac2003}, which minimizes the worst
weighted deviation among all normalized objectives.

\subsubsection*{Properties}
\begin{itemize}
  \item \emph{Full Pareto coverage}: Unlike the weighted-sum methods, the
    Chebyshev scalarization can recover Pareto optimal solutions on non-convex
    segments of the front~\cite{das1997, miettinen1998}. For any
    Pareto-efficient point $\mathbf{f}^*$, there exists an $\alpha$ such that
    $J_3$ achieves it.
  \item \emph{Non-differentiability}: The $\max$ operator introduces ridge lines
    where $(1-\alpha)\tilde{f}_1 = \alpha\tilde{f}_2$; gradient-based methods
    require subgradient extensions.
  \item \emph{Augmented form}: The regularized variant
    $J_3^{+} = J_3 + \rho(\tilde{f}_1 + \tilde{f}_2)$ with $\rho \ll 1$
    restores differentiability while maintaining full Pareto
    coverage~\cite{messac2003}.
  \item \emph{Discontinuous optimal trajectory}: As $\alpha$ varies, the
    argmin of $J_3$ may jump discontinuously between Pareto-efficient regions,
    making sensitivity analysis more challenging than for $J_1$.
\end{itemize}

\subsection{\texorpdfstring{$J_4$}{J4} — Euclidean Distance to Ideal Point}

\begin{equation}
  J_4(\rRx) = \sqrt{\tilde{f}_1^{\,2} + \tilde{f}_2^{\,2}}
  \label{eq:J4}
\end{equation}

$J_4$ minimizes the Euclidean distance from the normalized objective vector to
the ideal (utopia) point $(0, 0)$ in the normalized objective space. No
trade-off parameter is required.

\subsubsection*{Properties}
\begin{itemize}
  \item \emph{Parameter-free}: produces a single canonical compromise solution,
    useful when no preference information is available.
  \item \emph{Differentiability}: $J_4$ is smooth almost everywhere (except at
    the origin, which is infeasible by definition).
  \item \emph{Single Pareto point}: the minimizer is the knee-point of the
    normalized Pareto front, i.e., the solution closest to utopia. All other
    trade-offs require a different $J$.
  \item \emph{Reference-point sensitivity}: the choice of reference point
    (utopia, anti-utopia, or a design-specified aspiration point) substantially
    changes the minimizer, requiring calibration with domain
    knowledge~\cite{deb2001}.
\end{itemize}

\subsection{\texorpdfstring{$J_5$}{J5} — Logarithmic Sum with Brewster Penalty}

\begin{equation}
  J_5(\rRx;\,\alpha,\,\gamma)
  = (1-\alpha)\,f_1 + \alpha\,f_2
    + \gamma\,\bigl(\theta_R(\rRx) - \thetaB\bigr)^2
  \label{eq:J5}
\end{equation}

where $\theta_R(\rRx)$ is the angle of incidence of the Rx$\to$target ray at
the air--soil interface, and $\thetaB$ is given by~\eqref{eq:thetaB}.

\subsubsection*{Properties}
\begin{itemize}
  \item \emph{Physical motivation}: The quadratic penalty acts as a soft
    constraint driving the optimal Rx position toward the Brewster locus
    $\{\rRx : \theta_R = \thetaB\}$; for $\gamma \to \infty$ the solution
    converges exactly to this curve.
  \item \emph{Differentiability}: $J_5$ inherits the smoothness of $J_1$ plus
    the quadratic term, which is differentiable everywhere.
  \item \emph{Recovery}: For $\gamma = 0$, $J_5 \equiv J_1$.
  \item \emph{Redundancy concern}: $J_1$ already captures the Brewster effect
    through $f_1 = -\log_{10} P_r$, since $T_2(\theta_R)$ exhibits a sharp
    maximum at $\thetaB$ that dominates the landscape of $f_1$. Adding an
    explicit angular penalty may be redundant for the two-layer model.
  \item \emph{Calibration overhead}: introduces a third hyperparameter
    $\gamma \geq 0$, adding complexity relative to $J_1$.
  \item \emph{Hard-constraint alternative}: The $\varepsilon$-constrained
    formulation
    \begin{equation}
      \min_{\rRx}\,f_2(\rRx)\quad
      \text{s.t.}\quad f_1(\rRx) \leq f_1^{\max}
      \label{eq:eps_constrained}
    \end{equation}
    enforces the power requirement as a strict inequality rather than a soft
    penalty, which is preferred when a minimum SNR must be \emph{guaranteed}.
\end{itemize}

% ============================================================================
\section{Simulation Results}
\label{sec:results}
% ============================================================================

All simulations are performed analytically on a $270\times270$-point
Cartesian grid at $z_{\Rx} = 100\,\mathrm{m}$. Refracted ray parameters
are computed via vectorized bisection of Fermat's principle~\eqref{eq:fermat}
at each grid point. The generated figures referenced below are produced by the
script \texttt{utils/generate\_figures\_cost\_analysis.py}.

\subsection{Optimization Landscapes}

Fig.~\ref{fig:landscapes} shows the landscapes of the elementary objectives
($f_1$, $f_2$) and the four parametric cost functions ($J_1$, $J_2$, $J_3$,
$J_5$) in the $(x_{\Rx}, y_{\Rx})$ plane for $\alpha = 0.30$ and
$\gamma = 5\,\mathrm{rad}^{-2}$.

\begin{figure}[!t]
  \centering
  \includegraphics[width=\columnwidth]{../utils/fig_cf01_landscapes.pdf}
  \caption{Cost function landscapes in the $(\mathit{x}_{\mathrm{Rx}},
    \mathit{y}_{\mathrm{Rx}})$ plane for $\alpha = 0.30$.
    Gold star: Tx position projected onto the search plane.
    White cross: target projection on the surface.
    White circle: optimal Rx position (minimum of each function).}
  \label{fig:landscapes}
\end{figure}

Four principal features are observed:

\begin{enumerate}
  \item The landscape of $f_1$ (power loss) exhibits a narrow high-transmittance
    ring around the target at the radial distance $\rho_B \approx 200\,\mathrm{m}$
    corresponding to the Brewster angle (${\approx}5\gdeg$ angular width).
  \item The landscape of $f_2$ (resolution) shows that lateral resolution
    $\dxy$ improves monotonically as $|\Delta\phi|$ increases, with the global
    minimum at $\Delta\phi = 180\gdeg$ (Rx diametrically opposite to Tx).
  \item The minima of $J_1$ and $J_5$ are spatially proximal, with $J_5$
    applying a mild shift toward the Brewster locus for $\gamma = 5$.
  \item $J_2$ (linear normalized) exhibits extended plateau regions away from
    the Brewster ring, while $J_3$ (Chebyshev) produces a more compact minimum
    region with characteristic ridges from the $\max$ operator.
\end{enumerate}

\subsection{Pareto Fronts}

Fig.~\ref{fig:pareto} shows the Pareto fronts traced by varying $\alpha
\in [0,1]$ with a step of $0.01$ and computing the optimal Rx for each
function.

\begin{figure}[!t]
  \centering
  \includegraphics[width=0.9\columnwidth]{../utils/fig_cf02_pareto_fronts.pdf}
  \caption{Pareto fronts in the objective space $(f_1, f_2)$ for the five cost
    functions. Each symbol is the optimal Rx for one value of $\alpha$.
    Black star: ideal (utopia) point. Diamond ($J_4$): nearest-to-ideal point.}
  \label{fig:pareto}
\end{figure}

Three key observations follow:
\begin{enumerate}
  \item $J_1$ traces a smooth, nearly convex front. The high-curvature region
    near $\alpha \in [0.20, 0.40]$ is the Pareto knee, where small changes in
    $\alpha$ produce large movements of the optimal Rx.
  \item $J_3$ (Chebyshev) covers a strictly larger subset of the objective
    space than $J_1$ for $\alpha \in [0.3, 0.7]$, confirming the theoretical
    result that Chebyshev scalarization reaches non-convex
    segments~\cite{das1997}.
  \item $J_4$ produces a single point near the Pareto knee of $J_1$, consistent
    with the knee being the nearest point to the utopia in normalized space.
\end{enumerate}

\subsection{Sensitivity to Trade-Off Parameter}

Fig.~\ref{fig:alpha} plots the optimal Rx polar coordinates
$(\rho_{\Rx}^*,\,\phi_{\Rx}^*)$ as functions of $\alpha$ for $J_1$, $J_2$,
and $J_3$.

\begin{figure}[!t]
  \centering
  \includegraphics[width=\columnwidth]{../utils/fig_cf03_alpha_sensitivity.pdf}
  \caption{(a) Optimal receiver radial distance $\rho_{\mathrm{Rx}}^*$ versus
    $\alpha$. (b) Optimal receiver azimuth $\phi_{\mathrm{Rx}}^*$ versus
    $\alpha$. Compared for $J_1$, $J_2$, $J_3$.}
  \label{fig:alpha}
\end{figure}

For $J_1$, both $\rho_{\Rx}^*$ and $\phi_{\Rx}^*$ vary continuously and
monotonically with $\alpha$: for small $\alpha$ the Rx converges to the
Brewster ring ($\rho \approx 200\,\mathrm{m}$); for large $\alpha$ it moves
to maximize $|\Delta\phi|$. $J_3$ exhibits abrupt jumps at specific $\alpha$
values, a structural property of the Chebyshev scalarization arising from
branch switching in the $\max$ operator~\cite{miettinen1998}. The transition
region $\alpha \in [0.25, 0.35]$ shows the highest variability for $J_1$,
indicating that the choice of $\alpha^*$ within this interval requires careful
calibration.

\subsection{Gradient Conditioning}

Fig.~\ref{fig:grad} shows the normalized gradient norm $\norm{\nabla J}/
\max\norm{\nabla J}$ along the profile $y_{\Rx} = 0$ (a) and the gradient
norm map of $J_1$ with overlaid isolines (b).

\begin{figure}[!t]
  \centering
  \includegraphics[width=\columnwidth]{../utils/fig_cf04_gradient_norm.pdf}
  \caption{(a) Normalized gradient norm $\|\nabla J\|$ along the profile
    $\mathit{y}_{\mathrm{Rx}} = 0$ for $\alpha = 0.30$.
    (b) Gradient norm map of $J_1$ with $J_1$ isolines superimposed.}
  \label{fig:grad}
\end{figure}

$J_1$ maintains appreciable gradient throughout the domain, including regions
far from the minimum. $J_2$ develops near-zero gradient over large plateau
regions, degrading convergence of gradient descent from arbitrary starting
points. $J_3$ has high but discontinuous gradient at the ridge lines of the
$\max$ operator. The condition number $\kappa(\nabla^2 J_1) =
\lambda_{\max}/\lambda_{\min}$ of the Hessian of $J_1$ is more favorable than
that of $J_2$, resulting in faster convergence for quasi-Newton
methods~\cite{nocedal2006}.

\subsection{Effect of the Brewster Penalty}

Fig.~\ref{fig:brewster} characterizes the dependence of the optimal solution
on the penalty coefficient $\gamma$ in $J_5$.

\begin{figure}[!t]
  \centering
  \includegraphics[width=\columnwidth]{../utils/fig_cf05_brewster_penalty.pdf}
  \caption{(a) Incidence angle $\theta_R^*$ at the optimal Rx versus $\gamma$,
    with the asymptote $\thetaB$ (dashed). (b) Objective values $f_1^*$ and
    $f_2^*$ at the optimal Rx versus $\gamma$.}
  \label{fig:brewster}
\end{figure}

The optimal incidence angle converges rapidly to $\thetaB$ for
$\gamma \lesssim 5\,\mathrm{rad}^{-2}$ and saturates for
$\gamma > 10\,\mathrm{rad}^{-2}$. Panel~(b) confirms the expected trade-off:
increasing $\gamma$ reduces $f_1$ (higher power) but increases $f_2$ (coarser
resolution). The range $\gamma \in [1, 5]\,\mathrm{rad}^{-2}$ yields
approximately $85\%$ convergence toward $\thetaB$ with less than $15\%$
degradation in resolution relative to $\gamma = 0$.

% ============================================================================
\section{Discussion}
\label{sec:discussion}
% ============================================================================

\subsection{Why Logarithmic Scalarization Is the Preferred Formulation}

The fundamental advantage of $J_1$ lies in the logarithmic transformation of
$P_r$. The received power in~\eqref{eq:Pr} spans three to four orders of
magnitude across the search domain due to the combined effect of $R_R^{-2}$
path loss and the sharp Brewster transmittance peak. Applying $\log_{10}$
compresses this range to $[14, 18]$, placing $f_1$ and $f_2$ in numerically
comparable scales without requiring domain-dependent normalization. This is a
non-trivial property: the commensurability of $J_1$ is \emph{intrinsic} to the
physical model, not imposed externally, which means it is preserved as the
system parameters change.

Furthermore, the logarithmic function is concave, guaranteeing that the
gradient of $f_1$ is largest precisely near the Brewster angle (where $P_r$
changes most rapidly) and smallest in the flat plateau regions far from
$\theta_B$. This concentrates optimization effort where the physical response
is most informative.

\subsection{Conditions Favoring Chebyshev Scalarization}

In the two-layer half-space model considered here, the Pareto front is nearly
convex, making $J_1$ sufficient for exploring most Pareto-efficient solutions.
However, if the medium extends to a multilayer stack (e.g., soil over a water
table or permafrost layer), the TM transmittance can exhibit multiple local
maxima as a function of angle~\cite{balanis2012}, inducing a non-convex Pareto
front. In that case, $J_3$ with the augmented regularization
$\rho = 10^{-3}$ is the recommended alternative: it provides full Pareto
coverage while recovering differentiability for gradient-based search.

\subsection{Epsilon-Constrained Formulation as a Design Tool}

When the minimum received power is a hard system requirement (e.g., a minimum
SNR threshold for detection), the penalty approach of $J_5$ is inferior to the
$\varepsilon$-constrained formulation~\eqref{eq:eps_constrained}. The latter
guarantees that the power constraint is satisfied exactly for any feasible
solution, rather than balancing it softly against the resolution objective.
Modern constrained optimization methods such as augmented Lagrangian or
interior-point methods can solve~\eqref{eq:eps_constrained} with the same
efficiency as the unconstrained scalarizations, provided the constraint gradient
is computable—which it is, since $\partial f_1 / \partial \rRx$ is available
analytically via the chain rule through the refracted-ray
geometry~\eqref{eq:fermat}.

\subsection{Generalization to Multilayer Media}

The formulation extends naturally to an $N$-layer medium. The TM
transmittance product $T_2(\rRx)$ in~\eqref{eq:Pr} is replaced by the total
power transmittance of the stratified stack computed via the transfer matrix
(ABCD) method~\cite{balanis2012}. The expressions for $\dxy$ and $\dz$ require
updating the group velocity and the cumulative refraction angle through all
layers, but the structure of $f_1$ and $f_2$ as logarithmic objectives remains
unchanged. The five cost functions $J_1$–$J_5$ therefore generalize without
modification; only the underlying physics model is replaced.

% ============================================================================
\section{Conclusion}
\label{sec:conclusion}
% ============================================================================

This paper has presented a systematic analysis of five scalarization cost
functions for the problem of optimal bistatic receiver placement in helical SAR
subsurface imaging. The main conclusions are:

\begin{enumerate}
  \item The \emph{logarithmic weighted-sum} cost function $J_1$ is the
    recommended choice for the two-layer (air--soil) bistatic SAR geometry. Its
    intrinsic dimensional commensurability, continuous differentiability, and
    unimodal landscape topology make it well-suited for efficient gradient-based
    or quasi-Newton optimization. The recommended trade-off weight
    $\alpha^* = 0.30$ provides balanced power-resolution performance at the
    knee of the Pareto curve.

  \item The \emph{Chebyshev scalarization} $J_3$ should be adopted when the
    propagation medium is multilayer and the resulting Pareto front is non-convex.
    The augmented form with $\rho = 10^{-3}$ restores differentiability while
    preserving full Pareto-front coverage.

  \item The \emph{linear normalized weighted sum} $J_2$ is not recommended for
    production use: its domain-dependent normalization and gradient plateau
    regions result in poor numerical conditioning that undermines convergence
    reliability.

  \item The \emph{Brewster penalty term} in $J_5$ provides a physically
    motivated design knob but is largely redundant with $f_1$ in the two-layer
    case. The $\varepsilon$-constrained formulation~\eqref{eq:eps_constrained}
    is preferred when a hard SNR floor must be guaranteed.

  \item All five cost functions generalize to multilayer subsurface models via
    transfer-matrix propagation without altering the mathematical structure of
    the scalarization.
\end{enumerate}

Future work will extend the analysis to wideband configurations with frequency-
dependent permittivity, include the effect of surface roughness on the TM
transmittance model, and study the sensitivity of the optimal Rx position to
uncertainties in the soil electrical parameters.

% ============================================================================
%  REFERENCES
% ============================================================================
\bibliographystyle{IEEEtran}
\begin{thebibliography}{99}

\bibitem{moreira2013}
A.~Moreira, P.~Prats-Iraola, M.~Younis, G.~Krieger, I.~Hajnsek,
and K.~P.~Papathanassiou,
``A tutorial on synthetic aperture radar,''
\emph{IEEE Geosci.\ Remote Sens.\ Mag.}, vol.~1, no.~1, pp.~6--43,
Mar.~2013.

\bibitem{cumming2005}
I.~G.~Cumming and F.~H.~Wong,
\emph{Digital Processing of Synthetic Aperture Radar Data: Algorithms and
Implementation}.
Norwood, MA, USA: Artech House, 2005.

\bibitem{daniels2004}
D.~J.~Daniels,
\emph{Ground Penetrating Radar}, 2nd~ed.
London, U.K.: IET, 2004.

\bibitem{yarovoy2007}
A.~G.~Yarovoy, T.~G.~Savelyev, P.~J.~Aubry, P.~E.~Lys, and L.~P.~Ligthart,
``UWB radar for human being detection,''
\emph{IEEE Trans.\ Aerosp.\ Electron.\ Syst.}, vol.~43, no.~1,
pp.~22--31, Jan.~2007.

\bibitem{willis1991}
N.~J.~Willis,
\emph{Bistatic Radar}.
Norwood, MA, USA: Artech House, 1991.

\bibitem{krieger2006}
G.~Krieger \emph{et al.},
``TanDEM-X: A satellite formation for high-resolution SAR
interferometry and digital elevation model generation,''
\emph{IEEE Trans.\ Geosci.\ Remote Sens.}, vol.~45, no.~11,
pp.~3317--3341, Nov.~2006.

\bibitem{reigber2000}
A.~Reigber and A.~Moreira,
``First demonstration of airborne SAR tomography using multibaseline
L-band data,''
\emph{IEEE Trans.\ Geosci.\ Remote Sens.}, vol.~38, no.~5,
pp.~2142--2152, Sep.~2000.

\bibitem{tebaldini2012}
S.~Tebaldini and F.~Rocca,
``Multibaseline polarimetric SAR tomography of a boreal forest at P-
and L-bands,''
\emph{IEEE Trans.\ Geosci.\ Remote Sens.}, vol.~50, no.~1,
pp.~232--246, Jan.~2012.

\bibitem{born1999}
M.~Born and E.~Wolf,
\emph{Principles of Optics: Electromagnetic Theory of Propagation,
Interference and Diffraction of Light}, 7th~ed.
Cambridge, U.K.: Cambridge Univ.\ Press, 1999.

\bibitem{balanis2012}
C.~A.~Balanis,
\emph{Advanced Engineering Electromagnetics}, 2nd~ed.
Hoboken, NJ, USA: Wiley, 2012.

\bibitem{ulaby2014}
F.~T.~Ulaby, D.~G.~Long, W.~J.~Blackwell, C.~Elachi, A.~K.~Fung,
C.~Ruf, K.~Sarabandi, H.~A.~Zebker, and J.~Van~Zyl,
\emph{Microwave Radar and Radiometric Remote Sensing}.
Ann Arbor, MI, USA: Univ.\ Michigan Press, 2014.

\bibitem{soumekh1999}
M.~Soumekh,
\emph{Synthetic Aperture Radar Signal Processing with MATLAB Algorithms}.
New York, NY, USA: Wiley, 1999.

\bibitem{devaney1984}
A.~J.~Devaney,
``Geophysical diffraction tomography,''
\emph{IEEE Trans.\ Geosci.\ Remote Sens.}, vol.~GE-22, no.~1,
pp.~3--13, Jan.~1984.

\bibitem{miettinen1998}
K.~Miettinen,
\emph{Nonlinear Multiobjective Optimization}.
Boston, MA, USA: Kluwer Academic, 1998.

\bibitem{deb2001}
K.~Deb,
\emph{Multi-Objective Optimization Using Evolutionary Algorithms}.
Chichester, U.K.: Wiley, 2001.

\bibitem{das1997}
I.~Das and J.~E.~Dennis,
``A closer look at drawbacks of minimizing weighted sums of objectives
for Pareto set generation in multicriteria optimization problems,''
\emph{Struct.\ Optim.}, vol.~14, pp.~63--69, 1997.

\bibitem{messac2003}
A.~Messac, A.~Ismail-Yahaya, and C.~A.~Mattson,
``The normalized normal constraint method for generating the Pareto
frontier,''
\emph{Struct.\ Multidisc.\ Optim.}, vol.~25, no.~2, pp.~86--98,
Jul.~2003.

\bibitem{nocedal2006}
J.~Nocedal and S.~J.~Wright,
\emph{Numerical Optimization}, 2nd~ed.
New York, NY, USA: Springer, 2006.

\bibitem{zadeh1963}
L.~A.~Zadeh,
``Optimality and non-scalar-valued performance criteria,''
\emph{IEEE Trans.\ Autom.\ Control}, vol.~8, no.~1, pp.~59--60,
Jan.~1963.

\end{thebibliography}

\end{document}
